% Conclusiones del corte documentado; correspondencia con los seis objetivos
% específicos de la introducción. Extensión prevista: no más de tres páginas.
\chapter{Conclusiones}
\label{ch:conclusiones}

El trabajo establece un marco metodológico para relacionar instrumentación distribuida, modelos físicos y aprendizaje automático con el control de temperatura y humedad de un secador solar tipo invernadero. Respecto del objetivo general de la Sección~\ref{intro:objetivos}, se concretaron procedimientos para preparar datos, verificar balances y contrastar modelos. La validación del caso Coffee Solutions queda cerrada sobre el conjunto experimental de referencia y las tareas evaluadas en esta investigación. Las conclusiones se organizan según los seis objetivos específicos y distinguen este resultado del análisis de las estrategias de control.

\begin{enumerate}
\item \textbf{Caracterización del secador.} El análisis experimental describió la evolución de temperatura y humedad relativa en $36$ posiciones durante tres días, identificó diferencias entre superficies y cuantificó la dispersión espacial. Se establecieron temperatura, razón de humedad, captación solar, almacenamiento térmico y renovación exterior como elementos del modelo. El calentamiento con contenido de vapor constante del caso CFD mostró un comportamiento compatible con la psicrometría: disminuir humedad relativa no demuestra extracción de agua. La caracterización delimita las variables y el soporte del contraste realizado.

\item \textbf{Desarrollo de modelos determinísticos y basados en datos.} Se dispuso de representaciones complementarias de cavidad, CFD tridimensional, red reducida, ensambles y PINN. Los balances de la red y el contraste de integradores respaldaron su coherencia numérica; el refinamiento espacial oscilatorio de la cavidad delimitó la estimación de incertidumbre de discretización. Sobre el bloque experimental de prueba, a $1200$ s los ensambles redujeron el RMSE frente a persistencia aproximadamente $53{,}21\%$ en temperatura DS, $49{,}77\%$ en DHT y $42{,}65\%$ en humedad. Las representaciones de cuatro y seis estados conservan sus configuraciones respectivas al interpretar estos resultados.

\item \textbf{Plataforma de adquisición y preparación de datos.} Se documentaron arquitectura, identidad espacial, sincronización y procesamiento causal. La auditoría de $466\,560$ registros mostró integridad del conjunto experimental y permitió recuperar $51\,836$ antecedentes por canal asociados a la pared posterior. Esta corrección evitó una exclusión espacial sistemática durante el entrenamiento. La cobertura nominal del $100\%$ y la separación de $20$ s sustentan la preparación de las ventanas utilizadas en los experimentos; su interpretación operativa distingue cobertura del archivo y latencia de comunicaciones.

\item \textbf{Validación y análisis estadístico.} Se completó la validación del caso de estudio sobre el conjunto experimental de referencia mediante separación cronológica, métricas por variable y horizonte y comparación emparejada frente a persistencia. Los intervalos exploratorios conservaron el signo favorable de los ensambles, mientras que el árbol empeoró el canal DS a corto plazo. El objetivo de validación queda satisfecho para las tareas, posiciones y horizontes evaluados. La dependencia temporal y las diferencias entre configuraciones delimitan el alcance estadístico y evitan extender las conclusiones a regímenes no contrastados.

\item \textbf{Formulación y evaluación de control inteligente.} Se formularon PID, MRAC y MPC bajo acciones, restricciones y escenarios comunes. La interpretación física establece que la renovación higrométrica depende de la diferencia de razón de humedad con el exterior y que la recirculación actúa sobre mezcla e intercambio térmico. Los pronósticos evaluados carecen del contraste de sensibilidad a comandos necesario para comparar acciones candidatas. Por ello, la formulación está desarrollada y su evaluación comparativa permanece pendiente; las latencias de inferencia no permiten atribuir seguimiento, ahorro energético o factibilidad al ciclo completo.

\item \textbf{Alcance y aplicación a otros secadores.} La estructura conservativa, los operadores de observación y la separación temporal de fuentes constituyen los componentes con mayor potencial de adaptación. Capacidades, conductancias, propiedades ópticas, caudales y modelos de datos requieren caracterización local. El recinto vacío proporciona una base para estudiar el microclima, pero representar secado exige incorporar y validar almacenamiento, resistencia al flujo y evaporación del producto. Los resultados actuales no determinan tiempo de secado, calidad del café ni una configuración mínima de sensores y actuadores.
\end{enumerate}

La hipótesis de la Sección~\ref{intro:hipotesis} recibe respaldo de la integración de modelos, los balances verificados y la validación con el conjunto experimental. La contribución alcanzada es una metodología auditable para representar y predecir las variables evaluadas en Coffee Solutions, con dominio de aplicación explícito. Su aplicación al diseño de control se apoya en esta validación y exige examinar además la respuesta a acciones y el ciclo de regulación, conforme a la Sección~\ref{dis:control}.
